\documentclass[10pt,a4paper]{amsart}
\usepackage[margin=19mm]{geometry}
\usepackage{amsmath,amssymb,mathtools}
\usepackage[scr=rsfs,cal=euler]{mathalfa}
\usepackage{xfrac}
\usepackage[unicode,hidelinks,bookmarks=false]{hyperref}
\newcommand{\ie}{i.e.\ }
\newcommand{\cf}{cf.\ }
\newcommand{\resp}{respectively\ }
\newcommand{\Subsection}{\subsection}
\providecommand{\nobreakdash}{\nobreak\hbox{-}}
\setlength{\emergencystretch}{2em}
\multlinegap0pt
\usepackage{bm}
\usepackage{bbm}
\usepackage{dsfont}
\usepackage{xspace}
\usepackage{cancel}
\usepackage{mathtools}
\usepackage[shortlabels]{enumitem}
\usepackage{amsthm}
\makeatletter
\@ifundefined{theorem}{\newtheorem{theorem}{Theorem}}{}
\@ifundefined{lemma}{\newtheorem{lemma}[theorem]{Lemma}}{}
\@ifundefined{proposition}{\newtheorem{proposition}[theorem]{Proposition}}{}
\makeatother
\makeatletter
\newcommand{\id}{\operatorname{id}}
\expandafter\ifx\csname manualtheorem\endcsname\relax
\newenvironment{manualtheorem}[1]{%
  \renewcommand\themanualtheoreminner{#1}%
  \manualtheoreminner
}{\endmanualtheoreminner}
\fi
\makeatother
\title[Generalized Blaschke--Santal\'o-type inequalities, without s]{Generalized Blaschke--Santal\'o-type inequalities, without symmetry restrictions}
\author{Thomas A. Courtade, Edric Wang}
\date{Source: arXiv:2509.08998v1 (CC BY 4.0)}
\begin{document}
\maketitle

\noindent\emph{Source and licence.} Generalized Blaschke--Santal\'o-type inequalities, without symmetry restrictions. Version arXiv:2509.08998v1 (CC BY 4.0), distributed under the Creative Commons Attribution 4.0 International licence, as recorded in the arXiv licence metadata for this exact version and shown on the abstract page https://arxiv.org/abs/2509.08998v1. Full rights notice: licenses/arxiv-2509.08998-CC-BY-4.0.txt.\par\smallskip

\noindent\textit{Editorial scope.} Counted unit: the Theorem whose source label is \texttt{thm:EntropicDuality} in the article (source0.tex lines 399--416, source label \texttt{thm:EntropicDuality}), delivered together with the article's own complete original proof of it (source0.tex lines 449--489, one \texttt{proof} environment of 41 lines). No mathematics is written, repaired or re-derived: statement and proof are the article's own text line for line. The proof is detached from the statement in the source: the article states the result at lines 399--416 and prints its proof at lines 449--489, and the proof environment's own header names the statement, so the association is the article's own. Required context retained uncounted: prop:Qblocks (lines 274--279); lem:PropertiesDFunc (lines 419--420). Every definition, display and label of those lines is retained unchanged. Omitted from this package and not redistributed: 1--273, 280--398, 417--418, 421--448, 490--685. Nothing inside a retained excerpt is omitted or altered: every retained body is delivered exactly as the article's lines are written, so no omission receipt of any kind accompanies this package. Citations: the retained excerpts carry no citation command at all, so no bibliography entry of the article is transcribed and no reference is redistributed. Reference adaptation: none was needed and none was made; every reference inside the delivered unit resolves inside it, so no unresolved reference or citation is produced. Printed slips kept literal: no printed word, symbol, hypothesis or number of the counted statement or of its proof was altered. Layout and class: the article's own class is \texttt{article} with packages including hyperref, amsthm, amsmath, amssymb, enumerate, color, fullpage; the wrapper is the lane's pinned \texttt{amsart} editorial wrapper at 10pt on A4, which restates the theorem-like environments the retained text needs and adds the article's own macro definitions that the retained text needs (\texttt{\textbackslash id}), with extra wrapper packages (bm, bbm, dsfont, xspace, cancel, mathtools, enumitem). The delivered numbering is the unit's own. Assets: no retained span contains a figure environment, a graphics-inclusion command, a PDF-page inclusion, a table or a TikZ picture; no image file is copied into this package, the package declares no asset dependency and no image file is redistributed with it. Licence: version arXiv:2509.08998v1 of this article is distributed under the Creative Commons Attribution 4.0 International licence, as recorded in the arXiv licence metadata for this exact version and shown on the abstract page https://arxiv.org/abs/2509.08998v1. A full rights notice is delivered at licenses/arxiv-2509.08998-CC-BY-4.0.txt. Characters: every retained excerpt is pure ASCII, so the delivered engine, XeLaTeX with its default Latin Modern text font, reports no missing character.
\begin{proposition}\label{prop:Qblocks}The following hold:
\begin{enumerate}[(i)]
\item If $\pi_i Q \pi^*_i \not\geq 0$ for some $i\in \{1,\dots, n\}$, then $D_g(\mathbf{c},Q) = +\infty$.  
\item If $\pi_i Q \pi^*_i \geq 0$ for every $i\in \{1,\dots, n\}$, then $D(\mathbf{c},Q+\delta \id_E) <\infty$ for all $\delta>0$. 
\end{enumerate}
\end{proposition}

\begin{lemma}\label{lem:PropertiesDFunc} If the sharp constant $D$ such that \eqref{IntegralIneq} holds for all admissible functions satisfying \eqref{eq:pointwiseInequality} is finite,  then $\pi_i Q \pi^*_i \geq 0$ for each $i=1,\dots, n$.  
\end{lemma}

\begin{theorem}\label{thm:EntropicDuality}
Let $D \in \mathbb{R}\cup \{+\infty\}$. The following statements are equivalent.  
\begin{enumerate}[(A)]
\item If  even Borel functions $f_i: E_i \to \mathbb{R} \cup\{+\infty\}$   satisfy $e^{-f_i}\in L^1(E_i)$ and 
\begin{align}
\sum_{i=1}^n c_i f_i(x_i) \geq  \langle x, Q x\rangle, ~~~~ \forall x=(x_1, \dots, x_n) \in E, \label{eq:pointwiseInequality}
\end{align}
then 
\begin{align}
 \prod_{i=1}^n \left( \int_{E_i}e^{-f_i} d x_i\right)^{c_i} \leq e^{D}.  \label{IntegralIneq}
\end{align}
\item If $\mu_i \in \mathcal{P}(E_i)$, $i=1, \dots, n$, are symmetric then 
\begin{align}
  \sum_{i=1}^n c_i h(\mu_i) \leq \sup_{\mu \in \pi(\mu_1, \dots, \mu_n)}  \int_E \langle x, Q x\rangle d\mu(x) + D. \label{eq:entropyIneq}
\end{align}
\end{enumerate}
Moreover, log-concave $(\mu_i)_{i=1}^n$ saturate  \eqref{eq:entropyIneq} for the sharp constant $D= D_s(\mathbf{c},Q)$.
\end{theorem}

\begin{proof}[Proof of Theorem \ref{thm:EntropicDuality}]
$(B) \Rightarrow (A)$. Let $(f_i)_{i=1}^n$ be even Borel functions satisfying \eqref{eq:pointwiseInequality}.  Take  (symmetric) $\mu_i$ to have density proportional to $e^{-{f_i}} $. Assume for now that each $\mu_i$ has compact support (this is always possible by setting $f_i = +\infty$ outside the support of $\mu_i$, which preserves \eqref{eq:pointwiseInequality}).  Due to the assumption of compact support, $\mu_i$ has finite second moments, and therefore well-defined entropy $h(\mu_i) < \infty$.    In particular, 
$$
h(\mu_i) =   \log \left( \int e^{-f_i} d x_i \right)  + \int  f_i d\mu_i.
$$
By \eqref{eq:pointwiseInequality}, one may see that each $f_i$ is bounded from below on its (compact) essential support.  Since $h(\mu_i)<\infty$,   it follows that $f_i  \in L^1(\mu_i)$. 
 Thus, using definitions and (B), we have 
\begin{align*}
\sum_{i=1}^n c_i \left( \log \left( \int e^{-f_i} d x_i \right)  + \int  f_i d\mu_i  \right)&=  \sum_{i=1}^n c_i h(\mu_i) \\
&\leq  \max_{\mu \in \pi(\mu_1, \dots, \mu_n)} \int \langle x, Q x\rangle d\mu  +D\\
&\leq   \sum_{i=1}^nc_i \int f_i  d\mu_i  +D.
\end{align*}
Cancelling the $\int f_i  d\mu_i$ terms on both sides, we conclude (A) holds.  The assumption that $\mu_i$ is compactly supported can be removed using monotone convergence.

 

We now show $(A) \Rightarrow (B)$.    We can assume the sharp constant $D$ in \eqref{IntegralIneq} is finite, else the claim is trivially true.  Therefore, by Lemma \ref{lem:PropertiesDFunc}, we may assume that $\pi_i Q \pi_i^* \geq 0$ for each $i=1,\dots, n$.  

Fix symmetric measures $\mu_i \in \mathcal{P}(E_i)$, and assume without loss of generality that $h(\mu_i) > -\infty$.   
By the multimarginal Kantorovich duality, for any $\epsilon>0$, there exist even $f_i\in L^1(\mu_i)$ satisfying \eqref{eq:pointwiseInequality} and 
\begin{align}
\sum_{i=1}^n c_i \int f_i(x_i) d\mu_i  <   \max_{\mu \in \pi(\mu_1, \dots, \mu_n)} \int \langle x, Q x\rangle d\mu + \epsilon. \label{maxCouplingOrig}
\end{align}
Note that \eqref{eq:pointwiseInequality} implies
$$
c_i f_i(x_i) \geq \sup_{x_j \in E_j, j\neq i} \left\{ \langle x, Q x \rangle - \sum_{j=1, i\neq j}^n c_i f_i(x_i) \right\}.
$$
Since $\pi_i Q \pi_i^* \geq 0$, the RHS is a pointwise supremum of even functions that are convex in $x_i$, and is therefore  even and convex in $x_i$.  Both \eqref{eq:pointwiseInequality} and  inequality \eqref{maxCouplingOrig} are preserved by replacing $c_i f_i$ with this even convex function.  Repeating this same argument for each $i=1, \dots, n$ shows that we may assume without loss of generality that each $f_i$ is even and convex.  In fact, we may add $\epsilon' |x_i|^2$ to each $f_i$; this obviously preserves \eqref{eq:pointwiseInequality} for any $\epsilon'>0$; and, for $\epsilon'>0$ sufficiently small, \eqref{maxCouplingOrig} will also be preserved since the second moments $(\mu_i)_{i=1}^n$ are finite.  Therefore, we may assume without loss of generality that each $f_i$ is uniformly convex.  In particular, each $e^{-f_i} \in L^1(E_1)$.  



Now, starting with the variational representation for entropy and   (A), we have 
\begin{align*}
\sum_{i=1}^n c_i h(\mu_i)  &\leq \sum_{i=1}^n c_i \left( \int f_i d\mu_i +  \log\left( \int_{E_i} e^{-f_i} dx_i \right)   \right) 
\leq   \max_{\mu \in \pi(\mu_1, \dots, \mu_n)} \int \langle x, Q x\rangle d\mu + \epsilon + D.
\end{align*}
We conclude (A)$\Rightarrow$(B).  


If $\pi_iQ \pi^*_i \not\geq 0$ for some $i$, then Proposition \ref{prop:Qblocks} ensures that $D_s(\mathbf{c},Q)= +\infty$, and it suffices to consider  Gaussian (and therefore log-concave) measures to achieve  saturation.  On the other hand, if $\pi_iQ \pi^*_i  \geq 0$ for all $i=1, \dots, n$, then the argument above shows that it suffices to consider even convex $(f_i)_{i=1}^n$ satisfying \eqref{eq:pointwiseInequality} to saturate \eqref{IntegralIneq} for sharp constant $D$.  Since sharp constants in (A) and (B) are now established to be the same, returning to the argument in the proof of (B)$\Rightarrow$(A) shows that it suffices to consider log-concave $(\mu_i)_{i=1}^n$ to saturate \eqref{eq:entropyIneq}.
\end{proof}

\end{document}
